\section{Analysis of Agent Ranking and Effectiveness}
\label{sec:rank_analysis}

To validate our neuro-symbolic approach, we analyzed the correlation between the agent priority (\textit{Ranking}) and the actual extraction of requirements (\textit{Effectiveness}). The ranking is derived from the \texttt{original\_weights} assigned by the system, while effectiveness is defined as the presence of a non-zero \texttt{effective\_weight} in the final query composition.

\begin{table}[h]
\centering
\caption{Probability of Agent Contribution by Rank Position}
\label{tab:rank_effectiveness}
\begin{tabular}{lc}
\hline
\textbf{Agent Rank} & \textbf{Effectiveness Probability} \\
\hline
Rank 1 (Highest Priority) & 97.74\% \\
Rank 2 & 93.21\% \\
Rank 3 & 79.01\% \\
Rank 4 & 57.82\% \\
Rank 5 (Lowest Priority) & 33.95\% \\
\hline
\end{tabular}
\end{table}

The statistical analysis shows a strong correlation between the assigned rank and the probability of contributing requirements. We calculated a Spearman Correlation coefficient of $\rho = -0.5152$ ($p < 0.001$), confirming that agents in higher positions (Rank 1 and 2) are significantly more likely to provide relevant constraints. Notably, an agent in the first position is approximately 2.9 times more likely to contribute than an agent in the fifth position, empirically validating the system's ability to correctly prioritize agents based on the natural language query context.
